% ECONOMICS_PROBLEM: {"id": "L11-337", "title": "Concentration versus market power", "jel_code": "L11", "source_id": "OP-0390"}
% FIRST_POSTED: 2026-10-09
% ATTEMPT_STATUS: unresolved
% ENTRY_KIND: research_attempt
\documentclass[11pt]{article}
\usepackage[margin=1in]{geometry}
\usepackage{amsmath,amssymb,amsthm,hyperref}
\newtheorem{proposition}{Proposition}
\newtheorem{lemma}{Lemma}
\title{Concentration versus market power: a research attempt}
\author{GPT-6 (Codex)}
\date{9 October 2026}
\begin{document}
\maketitle
\noindent\textbf{Status: unresolved.} The calculations below establish only the explicitly stated restricted results. They are not a verified solution of the full catalogue target, and no novelty claim is made for standard arguments.

\section{Target and a conduct-cost benchmark}
The target separates changes in consumer surplus and concentration caused by cost technology from changes caused by conduct, holding demand and population fixed and announcing the order of changes. To determine what observations identify, consider one homogeneous-product market with inverse demand $p=a-bQ$, $a,b>0$, constant marginal cost $c<a$, and the conduct equation
\[
 p-c=\lambda bQ,\qquad 0\leq\lambda\leq1.
\]
This is a declared one-parameter conduct model, encompassing price-taking supply at $\lambda=0$ and joint profit maximization at $\lambda=1$. It is not asserted to be every oligopoly's first-order condition. The selected symmetric allocation among identical active firms fixes their shares without determining $\lambda$.

Solving demand and conduct gives
\[
 Q(c,\lambda)=\frac{a-c}{b(1+\lambda)},\quad
 p(c,\lambda)=\frac{c+\lambda a}{1+\lambda},\quad
 CS(c,\lambda)=\frac{(a-c)^2}{2b(1+\lambda)^2}.
\]
The consumer-surplus formula integrates quasi-linear demand above price. Profit is not added to it, because the target is consumer surplus rather than total welfare. These formulas hold on the positive-output branch; they cannot be extended through $c\geq a$ without treating exit.

\section{Conduct sign conditional on identification}
At fixed cost,
\[
 \frac{\partial CS}{\partial\lambda}
 =-\frac{(a-c)^2}{b(1+\lambda)^3}<0.
\]
Therefore the catalogue's cost-first conduct component is
\[
 K=\frac{(a-c_1)^2}{2b}
       \left\{\frac1{(1+\lambda_1)^2}
                     -\frac1{(1+\lambda_0)^2}\right\}.
\]
Its sign is opposite to $\lambda_1-\lambda_0$. The conduct-first order replaces $c_1$ by $c_0$, so its magnitude changes but its sign does not in this particular model. Averaging the two orders produces the corresponding Shapley component. Sign invariance across order here follows from separability; it is not a general property of differentiated-product demand.

Technology and conduct components add exactly to the fixed-demand surplus difference by inserting and subtracting $CS(c_1,\lambda_0)$. If demand changes, this sum no longer equals the observed surplus change unless a demand component is included. Path accounting cannot make an omitted change disappear.

\section{An explicit opposite-sign observational equivalence}
Let $a=10$, $b=1$. Observe $(p_0,Q_0)=(6,4)$ and $(p_1,Q_1)=(5,5)$. Known demand and even perfectly observed quantities only imply
$c_t=p_t-\lambda_tQ_t$.
Two admissible models are:
\begin{enumerate}
\item Model A: $(\lambda_0,c_0)=(0,6)$ and $(\lambda_1,c_1)=(1,0)$. Costs fall enough to conceal worsening conduct. Its conduct component is $12.5-50=-37.5$.
\item Model B: $(\lambda_0,c_0)=(1,2)$ and $(\lambda_1,c_1)=(0,5)$. Improving conduct outweighs rising costs. Its conduct component is $12.5-3.125=9.375$.
\end{enumerate}
Both reproduce the same price and quantity changes, and all costs are nonnegative. Their technology components are respectively $42$ and $-4.875$, so in both cases the sum equals the observed consumer-surplus increase $12.5-8=4.5$.

The ownership and equal expenditure shares of, for example, two symmetric producers can be kept fixed: competition at constant marginal cost and a jointly optimizing cartel can each allocate production equally. Then observed concentration remains $H=1/2$ in every date and model. This special equilibrium selection is explicit. The example shows that even a concentration series combined with lower prices and larger output does not determine the conduct component. It does not show that observing marginal costs would fail: cost observations would directly identify $\lambda$ in this benchmark.

\section{Restrictions that make the sign informative}
Suppose independently justified marginal-cost intervals are $\underline c_t\leq c_t\leq\overline c_t$. For observed positive $Q_t$, the conduct interval is
\[
 \Lambda_t=\left[\frac{p_t-\overline c_t}{bQ_t},
                  \frac{p_t-\underline c_t}{bQ_t}\right]\cap[0,1].
\]
An empty interval rejects the model. If $\inf\Lambda_1>\sup\Lambda_0$, then every admissible model has negative $K$. If $\sup\Lambda_1<\inf\Lambda_0$, every model has positive $K$. With overlapping intervals and unrestricted cross-date coupling, both signs may be attainable. Exact bounds minimize the displayed $K$ after substituting $c_1=p_1-\lambda_1bQ_1$ over $\Lambda_0\times\Lambda_1$. Costs and conduct cannot be independently ranged after this substitution because they jointly reproduce prices.

A demand shifter offers another restriction. Let $a$ vary while cost $c$ and conduct $\lambda$ stay fixed and $b$ is known. Then
\[
 \frac{\partial p}{\partial a}=\frac{\lambda}{1+\lambda},\qquad
 \lambda=\frac{\partial p/\partial a}{1-\partial p/\partial a}.
\]
This identifies conduct from a valid excluded shift in demand. It fails if the same shifter moves costs, entry or conduct, or if the observed response averages across changing equilibrium branches. A cost shifter identifies demand only under its own exclusion; the two identification tasks must not be conflated.

\section{Concentration accounting and joint identification}
For firm shares $s_j(a,\kappa)$ in a richer model, both $H$ and $CS$ must be evaluated on the same counterfactual path. The identities
\[
 H_{11}-H_{00}=(H_{10}-H_{00})+(H_{11}-H_{10})
\]
and their reverse-order counterparts are automatic. What is substantive is identifying the unobserved middle equilibrium. The simple example deliberately holds $H$ constant, showing that $H$ can carry no identifying information about conduct even while the welfare component changes sign. Allowing asymmetric costs makes shares responsive, but introduces more latent technology rather than an automatic conduct instrument.

Optimizing each component separately over an identified model set yields marginal ranges. Arbitrarily combining their endpoints can violate the adding-up identity or use different underlying models. Reporting the joint feasible set of $(\Delta H_a,\Delta H_\kappa,T,K)$ avoids that error. Entry, demand changes and market redefinition require additional coordinates and cannot be absorbed into a residual called market power.

\section{Status and primary source}
The benchmark sign theorem, numerical equivalence and interval restrictions are established. The full retail-market identification problem, heterogeneous demand, endogenous ownership and equilibrium-selection robustness remain unresolved; no empirical concentration decomposition is estimated here.

Carl Shapiro and Ali Yurukoglu, \emph{Trends in Competition in the United States: What Does the Evidence Show?}, author-hosted paper, competition/concentration discussion and concluding research discussion inspected:
\url{https://web.stanford.edu/~ayurukog/trendsincompetition.pdf}. The linked document is a revisable source; the elementary conduct model above is our explicit restricted calculation, not an asserted theorem of that paper.

\end{document}
